\chapter*{Executive Summary}
\addcontentsline{toc}{chapter}{Executive Summary}

University students spend much of their study time reading PDF textbooks and lecture slides, yet
reading alone gives little evidence of understanding, and the techniques that do build lasting
knowledge (structured notes, self-testing and spaced revision) take more time than most students can
give to every topic. Existing tools each cover one part of this: PDF readers display documents,
document-chat tools answer the student's questions, and flashcard applications schedule revision of
cards the student must write.

RecaLL is a local-first desktop application that joins these steps into one study path. A student
imports an academic PDF; RecaLL extracts its structure into a topic tree, writes Cornell notes for each
topic, questions the student through Socratic drills that test reasoning rather than recall, turns
detected knowledge gaps and misconceptions into targeted flashcards, and schedules their revision with
the Free Spaced Repetition Scheduler (FSRS-6). A Knowledge Hub answers questions across the student's
own documents with citations that are checked against the retrieved passages, quizzes test whole
documents, and a progress view shows what has been studied and what is due.

The system was developed incrementally between July and November 2026 by a team of four, led by Al
Rubayet. It uses a
layered client--server architecture inside a single desktop package: a Tauri~2 shell with a React~19 and
TypeScript interface, and a local FastAPI backend that stores all data in SQLite. AI requests go through
a provider layer, so the student can choose Gemini, OpenAI or Groq with a personal key (kept in a vault
encrypted with AES-256-GCM), or a local Ollama model that keeps all data offline. No model is trained;
behaviour is controlled by prompts, schema validation and citation checking.

The backend has an automated test suite of 223 pytest cases covering authentication, user isolation,
document processing, FSRS scheduling, retrieval, quizzes and security primitives, together with nine
Playwright end-to-end journeys and a GitHub Actions pipeline for testing and Windows release builds.
All 223 tests pass, with 50.5\% statement coverage of the application code. Measured on the team's own
documents, the backend needs at most 381~MB of memory, a Cornell note takes a median of 27.9~s, and
Knowledge Hub answers cited the relevant pages in every answerable test question.
Six students tested the application informally during development, and their five feature requests
were addressed in the current build.

RecaLL shows that the full path from a PDF to scheduled, evidence-based revision can run privately on a
student's own computer. Future work will add semantic embeddings by default and built-in OCR, merge
the two databases behind one authenticated API, add cloud sync, and provide installers for more
platforms.
